```latex
\documentclass{article}
\usepackage{amsmath,amssymb,amsthm}

\newtheorem{theorem}{Theorem}[section]
\newtheorem{lemma}[theorem]{Lemma}
\theoremstyle{remark}
\newtheorem{remark}[theorem]{Remark}

\newcommand{\M}{\mathcal{M}}
\newcommand{\C}{\mathbb{C}}
\newcommand{\D}{\Delta}
\newcommand{\G}{\Gamma}
\newcommand{\p}{\partial}
\newcommand{\barM}{\overline{\M}_g}

\begin{document}

\title{Non-varying boundary factors for deformations of totally geodesic subvarieties}
\author{}
\date{}
\maketitle

Let $\M_g$ denote the moduli space of genus $g$ Riemann surfaces, and let
$\barM$ be its Deligne--Mumford compactification.  We use the following
rigidity theorem of McMullen.

\begin{theorem}[McMullen, Rigidity of totally geodesic subvarieties]\label{thm:rigidity}
Let $Z\subset \M_g$ be a totally geodesic subvariety with respect to the
Teichm\"uller metric.  If $\mathcal Z\subset \M_g\times \D$ is a holomorphic
family of subvarieties of $\M_g$ deforming $Z$, i.e.\ $\mathcal Z_0=Z$, then
the family is trivial: after possibly shrinking $\D$, one has
$\mathcal Z_t=Z$ for all $t\in \D$.
\end{theorem}

\begin{proof}[Proof of Theorem~\ref{thm:rigidity}]
This is McMullen's rigidity theorem for totally geodesic subvarieties of
moduli space; see \cite{McMullen}, Theorem~1.2.
\end{proof}

We now prove the statement.

\begin{proof}
By hypothesis, $V_0\subset \M_g$ is a smooth totally geodesic subvariety of
complex dimension $2$, and
\[
\mathcal V\subset \M_g\times \D
\]
is a holomorphic family of subvarieties with fibers
$V_t=\psi(\phi^{-1}(t))$, where
$\phi:\mathcal V\to \D$ and $\psi:\mathcal V\to \M_g$ are the natural
projections.  Thus $\mathcal V$ is a holomorphic family of subvarieties of
$\M_g$ deforming $V_0$.

Applying McMullen's rigidity theorem (Theorem~\ref{thm:rigidity}) to
$Z=V_0$ and $\mathcal Z=\mathcal V$, we conclude that, after possibly
shrinking the base $\D$, the family is constant:
\[
V_t=V_0
\]
for all $t$ in the base.  Since $\mathcal V$ is a complex submanifold and
the base $\D$ is connected, equality on a neighbourhood of $0$ implies
equality over the whole base; hence we may take this to hold for all $t$ in
the original disc.

Now take the closure in $\barM$.  Since $V_t=V_0$ as subvarieties of
$\M_g$, their closures in $\barM$ coincide:
\[
\overline V_t=\overline V_0.
\]
Intersecting with the fixed boundary stratum $\D_\G$, we obtain
\[
\p V_t
:=\overline V_t\cap \D_\G
=\overline V_0\cap \D_\G
=:\p V_0.
\]
Therefore, for the projection $\pi_i:\D_\G\to \M_{g_i,n_i}$,
\[
\pi_i(\p V_t)=\pi_i(\p V_0).
\]
By assumption, $\pi_i(\p V_0)=\{Y_0\}$.  Hence
\[
\pi_i(\p V_t)=\{Y_0\}
\]
for all $t$, which proves that the family $\p V_t$ is non-varying on that
factor.
\end{proof}

\begin{remark}
The technical assumptions (i) and (ii) in the statement are not needed once
the full rigidity theorem for totally geodesic subvarieties is invoked.
They ensure that the boundary intersections form a well-behaved family, but
the rigidity theorem already forces the entire family $V_t$ to be constant.
\end{remark}

\begin{thebibliography}{1}
\bibitem{McMullen}
C.~T. McMullen,
\emph{Rigidity of Teichm\"uller curves},
Ann. of Math. (2) \textbf{170} (2009), no.~1, 267--302.
\end{thebibliography}

\end{document}
```